\emph{(i)} $P$ is real, symmetric and $P^{2} = 1$, so $X \mapsto PXP$ is an involutive algebra automorphism of $M_{3}(\C)$ commuting with the adjoint and the transpose ($P^{\dagger} = P^{\mathsf T} = P$); its fixed points are therefore closed under products, adjoints, transposes and inverses ($PX^{-1}P = (PXP)^{-1}$). Written out, $PXP = X$ means $X_{11} = X_{22}$, $X_{12} = X_{21}$, $X_{13} = X_{23}$ and $X_{31} = X_{32}$: five complex parameters (the check computes the fixed space of the nine linear conditions and finds dimension $5$; for symmetric matrices this is the four-dimensional space of Proposition \claimref{synthesis.residual-s2-fixes-120-family-form}(i), for antisymmetric ones the line $(e_{1} + e_{2})\wedge e_{3}$, and the $S_{3}$-invariant symmetric ones, $\alpha\mathbf 1 + \beta\J$ of Lemma \claimref{neutrino.invariant-symmetric-matrices}, are the special case). For such $X$, $Xa = (X_{11} - X_{12},\ X_{21} - X_{22},\ X_{31} - X_{32}) = (X_{11} - X_{12})\,a$, and $Xb_{1}$, $Xb_{2}$ lie in the even plane (their first two coordinates are equal), so in the basis $a, b_{1}, b_{2}$ every $S_{2}$-invariant matrix is block diagonal, $1 \oplus 2$. By Fact \claimref{standard-model.imported-clebsch-structure-seesaw}, $M_{u}, M_{d}, M_{e}, M_{\nu}$ are complex linear combinations of the Yukawa matrices $Y_{a}$, $f$, $h$ ($h$ antisymmetric, Fact \claimref{standard-model.imported-spinor-bilinears}) and $M_{R} = v_{R}f$, so if the Yukawas are $S_{2}$-invariant so are the five mass matrices, hence the Hermitian squares $MM^{\dagger}$ and the seesaw $m_{\nu} = -M_{\nu}M_{R}^{-1}M_{\nu}^{\mathsf T}$, by closure; $a$ is an eigenvector of each.

\emph{(ii)} Let $H$ be Hermitian with $Ha = \lambda a$ and $\lambda$ a simple eigenvalue. Any unitary $U$ with $U^{\dagger}HU$ diagonal has columns that are eigenvectors; the $\lambda$-eigenspace is $\C a$, so one column is $e^{i\phi}a/\sqrt2$. For the charged sector the left-handed mixing comes from $H_{e} = M_{e}M_{e}^{\dagger}$ (and $H_{u}$, $H_{d}$ for quarks), for the light neutrinos from the symmetric $m_{\nu}$ through $m_{\nu}m_{\nu}^{\dagger}$, whose eigenvectors are those of the Takagi factorisation when the singular values are simple. If $U_{1}$ and $U_{2}$ both have $e^{i\phi_{j}}\hat a$ as their column $k_{j}$, then $(U_{1}^{\dagger}U_{2})_{k_{1}k_{2}} = e^{i(\phi_{2} - \phi_{1})}\hat a^{\dagger}\hat a$ has modulus one, and unitarity forces the rest of row $k_{1}$ and column $k_{2}$ to vanish; the remaining $2 \times 2$ block is unitary and the Jarlskog invariant, a product over a $2 \times 2$ minor involving the vanishing entries, is zero. The labels are an extra input: the lemma places the unimodular entry in the row and column of $a$ in the two mass orderings, and the eigenvalue of $a$, $X_{11} - X_{12}$ in each matrix, need not be the smallest. At the $S_{2}$ stage $a$ is a null vector of $M_{D}$ and $M_{e}$ (Proposition \claimref{synthesis.residual-s2-fixes-120-family-form}(ii); re-verified in the check on the stored textures), so for the leptons $a$ is the lightest state in both sectors and the entry of modulus one is $U_{e1}$: $U_{e2} = U_{e3} = U_{\mu1} = U_{\tau1} = 0$, which in the standard parametrisation is $\theta_{12} = \theta_{13} = 0$. For the quarks the store does not fix the position of $a$ in the up and down orderings; if it is the lightest in both, $V_{us} = V_{ub} = V_{cd} = V_{td} = 0$, and otherwise the zeros fill another row and column (for instance $V_{ud} = 0$), which the record excludes as well. The $2 \times 2$ blocks are unconstrained. When $\lambda$ is degenerate (the light plane at leading order, where $m_{\nu}a = 0$ and the second light mass is not yet lifted) the eigenvector is not determined and neither is the mixing, which is the stored leading-order statement; the lemma bites once the second stage has lifted the degeneracies while keeping $S_{2}$.

\emph{(iii)} The $120$ of Proposition \claimref{synthesis.residual-s2-fixes-120-family-form}(ii) is the $S_{2}$-invariant antisymmetric line, $u'\wedge s = -3(e_{1} + e_{2})\wedge e_{3}$; it annihilates $a$ (checked) and by (i) it cannot move $a$ out of the eigenbasis of any mass matrix. So a second stage consisting of that $120$ alone, with $S_{2}$-invariant $10$s and $126$, is not a breaking of $S_{2}$ at all and yields the zeros of (ii): the second line of attack of the open problem, as phrased, cannot produce the Cabibbo angle or $\theta_{12}$, $\theta_{13}$. The family form of the $120$ is fixed by the residual $S_{2}$ only while that symmetry holds; a genuine second stage must add $S_{2}$-odd components to the charged-fermion couplings ($10$s, $126$ or $120$): by (i) an $S_{2}$-odd Majorana sector with $S_{2}$-invariant Dirac couplings (an $S_{2}$-odd $f$ entering only $M_{R} = v_{R}f$, that is $w_{u} = w_{d} = 0$ for the $126$ in the Dirac matrices) changes $m_{\nu}$ but none of $H_{e}$, $H_{u}$, $H_{d}$, and by Lemma \claimref{neutrino.breaking-relation} a breaking of the Majorana sector alone lifts no second light mass; with $w_{u}, w_{d} \neq 0$ an $S_{2}$-odd $f$ is itself an $S_{2}$-odd charged-fermion coupling. Proposition \claimref{synthesis.solar-angle-light-plane-sum-rule} excludes the one-type $S_{2}$-odd stages at the order of Lemma \claimref{neutrino.breaking-relation} and under Premise (H) ($m_{1}/m_{2} \le 0.01$), so under those two conditions what remains is an $S_{2}$-odd stage with both symmetric and antisymmetric parts in some sector: lopsided.

\emph{(iv)} The check builds $Y_{1}, Y_{2}, f$ (symmetric, $S_{2}$-invariant, complex rational) and $h = (1 + 2i)(e_{1} + e_{2})\wedge e_{3}$, the Clebsch combinations with the vevs $v_{u} = (3 + i, -1 + 2i)$, $v_{d} = (1 - i, 2)$, $w_{u} = \tfrac12 + \tfrac{i}{3}$, $w_{d} = -2 + i$, $120$ coefficients $(1 + i, 2i, \tfrac32, -1 + \tfrac{i}{2})$ and $v_{R} = 7$, the five mass matrices, the Hermitian squares and the seesaw (exact inverse of $M_{R}$), and verifies: commutation with $P$, $a$ an eigenvector with zero coupling to $b_{1}, b_{2}$, and the simplicity of its eigenvalue by $\det(B - \lambda) \neq 0$ on the even block $B$ for each Hermitian square and for $m_{\nu}$ ($H_{e}a = 205a$ pinned). The control adds $\tfrac13 e_{1}\wedge e_{3}$ to $h$: the $120$ is still antisymmetric but not $S_{2}$-invariant, and $a$ acquires a nonzero coupling to the even plane in $H_{e}$, $H_{d}$ and $m_{\nu}$ and is no longer an eigenvector of any of them. The stored textures at $d = 2$, $\delta = \tfrac27$, $\delta_{c} = -\tfrac27$, $\alpha = \tfrac15$, $M_{3} = 100$, $\varepsilon = \tfrac34$ (a breaking of the Majorana sector included) are $S_{2}$-invariant, $M_{D}a = M_{e}a = m_{\nu}a = 0$, $M_{R}a = \alpha a$, and the even block of $m_{\nu}$ is nonzero.
